\documentclass[a4paper,11pt]{article}
\usepackage[T2A]{fontenc}
\usepackage[utf8]{inputenc}
\usepackage[russian]{babel}
\usepackage{amsmath,amssymb}
\usepackage{geometry}
\geometry{left=23mm,right=22mm,top=23mm,bottom=27mm}
\usepackage{microtype}
\usepackage{parskip}
\usepackage{tabularx,booktabs,array}
\usepackage{xcolor}
\usepackage{tikz}
\usepackage{fancyhdr}
\usepackage{hyperref}
\definecolor{accent}{HTML}{008080}
\definecolor{darkaccent}{HTML}{004D4D}
\definecolor{textgray}{HTML}{333333}
\hypersetup{colorlinks=true,linkcolor=darkaccent,urlcolor=darkaccent}
\setlength{\parindent}{0pt}
\setlength{\parskip}{7pt}
\setlength{\emergencystretch}{2em}
\pagestyle{fancy}
\fancyhf{}
\fancyfoot[C]{\scriptsize Лёвин Артём Александрович эксклюзивно для Ярослава Гаврилова}
\fancyfoot[R]{\scriptsize\thepage}
\renewcommand{\headrulewidth}{0pt}
\renewcommand{\footrulewidth}{0pt}
\begin{document}
\hypersetup{pageanchor=false}
\begin{titlepage}
\thispagestyle{empty}
\begin{center}
\vspace*{24mm}
{\small\color{darkaccent} МАТЕМАТИКА\par}
\vspace{22mm}
{\LARGE Ревью домашней работы\par}
\vspace{16mm}
{\large Ярослав Гаврилов\par}
\vfill
{\normalsize 26 октября 2026 года\par}
\vspace{18mm}
{\small Лёвин Артём Александрович\par}
{\small эксклюзивно для Ярослава Гаврилова\par}
\vspace{17mm}
\begin{tikzpicture}[x=1cm,y=1cm]
\draw[accent,line width=0.9pt] (0,0) .. controls (1.35,0.65) and (2.6,-0.48) .. (4,0)
   .. controls (5.4,0.48) and (6.7,-0.45) .. (8,0);
\end{tikzpicture}
\end{center}
\end{titlepage}
\hypersetup{pageanchor=true}

\section*{Сводная таблица}
В таблицу включены только задания, для которых требуется завершить или исправить решение.
\vspace{6mm}
\begingroup
\small
\renewcommand{\arraystretch}{1.42}
\setlength{\tabcolsep}{3pt}
\begin{tabularx}{\dimexpr\linewidth-8pt\relax}{@{}p{10mm}p{24mm}X p{19mm}p{20mm}@{}}
\toprule
\textbf{№} & \textbf{Тема} & \textbf{Суть замечания} & \textbf{Ваш ответ} & \textbf{Верный ответ}\\
\midrule
\hyperref[task:5]{5} & Степенная зависимость & Решение не завершено: не получено уравнение для показателя степени. & Не указан & $a=3$\\
\bottomrule
\end{tabularx}
\endgroup

\clearpage
\thispagestyle{empty}
\vspace*{\fill}
\begin{center}
{\LARGE\color{darkaccent} Разбор ошибок}
\end{center}
\vspace*{\fill}
\clearpage

\phantomsection\label{task:5}
\section*{Задание № 5. Закон зависимости давления от объёма}

\textbf{Текст задания.} Давление газа и его объём связаны законом
\[
 pV^a=C.
\]
При уменьшении объёма втрое давление выросло в $27$ раз. Найдите $a$.

\textbf{Ваше решение.} В тетради присутствует зачёркнутый черновой фрагмент. Завершённой цепочки преобразований и окончательного результата нет. Зачёркнутую запись нельзя уверенно считать частью решения.

\textbf{Ваш ответ.} Не указан.

\textcolor{red}{\textbf{Ошибка:}} решение осталось незавершённым. Не составлено проверяемое уравнение, связывающее изменение давления с изменением объёма.

\textbf{Где именно возникает затруднение.} Последний достоверно правильный математический шаг в записи выделить нельзя: после номера задания нет завершённого вывода. Первый необходимый, но отсутствующий шаг --- обозначить исходные давление и объём и записать их значения после изменения.

\textbf{Почему этот шаг важен.} Равенство $pV^a=C$ выполняется как до изменения объёма, так и после него. Постоянная $C$ одна и та же. Слово «втрое» означает, что новый объём равен одной трети исходного, а «выросло в 27 раз» означает умножение исходного давления на $27$. Если эти изменения не подставить в один и тот же закон, показатель $a$ определить нельзя.

\textcolor{blue!60!black}{\textbf{Ключевая идея:}} сравнить два состояния газа при неизменном $C$ и сократить одинаковые ненулевые множители.

\textbf{Исправление.} Обозначим исходные величины $p_1>0$ и $V_1>0$. После изменения
\[
 p_2=27p_1,\qquad V_2=\frac{V_1}{3}.
\]
Теперь закон можно записать дважды и приравнять полученные выражения.

\textbf{Полное правильное решение.}

\emph{Шаг 1.} До изменения состояния газа выполняется
\[
 p_1V_1^a=C.
\]
\emph{Шаг 2.} После изменения давления и объёма выполняется
\[
 p_2V_2^a=C.
\]
Подставим $p_2=27p_1$ и $V_2=V_1/3$:
\[
 27p_1\left(\frac{V_1}{3}\right)^a=C.
\]
\emph{Шаг 3.} Поскольку обе записи равны $C$, приравниваем их:
\[
 27p_1\left(\frac{V_1}{3}\right)^a=p_1V_1^a.
\]
При $V_1>0$ используем свойство степени частного:
\[
 \frac{27p_1V_1^a}{3^a}=p_1V_1^a.
\]
\emph{Шаг 4.} Давление $p_1$ и объём $V_1$ положительны, поэтому множитель $p_1V_1^a$ не равен нулю. Делим на него обе части:
\[
 \frac{27}{3^a}=1,
 \qquad 3^a=27.
\]
\emph{Шаг 5.} Представляем $27$ степенью тройки:
\[
 3^a=3^3\quad\Longrightarrow\quad\boxed{a=3}.
\]

\textbf{Проверка.} При $a=3$ произведение после изменения равно
\[
 (27p_1)\left(\frac{V_1}{3}\right)^3
 =\frac{27p_1V_1^3}{27}=p_1V_1^3=C.
\]
Значит, требуемые изменения давления и объёма согласуются с законом.

\textcolor{teal!60!black}{\textbf{Итог:}} показатель степени равен $\boxed{3}$. При сравнении двух состояний сначала запишите, во сколько раз изменилась каждая величина, затем подставляйте эти изменения в исходную формулу.

\subsection*{Задача на закрепление}
Давление газа и объём связаны законом $pV^a=C$. Когда объём уменьшился вдвое, давление увеличилось в $16$ раз. Найдите $a$.

\end{document}
